\chapter{Manuale Tecnico}

\section{Introduzione}
Il sistema \textbf{Wiki UNIPR} è progettato per essere eseguito in un ambiente containerizzato tramite \textbf{Docker}. Questa architettura garantisce che il software funzioni in modo identico su qualsiasi macchina, isolando le dipendenze del Backend (PHP), del Database (MySQL) e del Frontend (Apache).

\section{Requisiti di Sistema}
Prima di procedere, assicurarsi che i seguenti strumenti siano installati e correttamente configurati:
\begin{itemize}
    \item \textbf{Docker Desktop} (Windows/macOS) o \textbf{Docker Engine} (Linux).
    \item \textbf{Docker Compose} (solitamente incluso nelle versioni recenti di Docker).
    \item Un terminale (Bash, PowerShell o WSL).
\end{itemize}

\section{Inizializzazione del Database}
Il database viene configurato automaticamente durante il primo avvio del container DB. 
\subsection{Logica di Popolamento}
Docker esegue gli script SQL presenti nella cartella \texttt{./services/database/configs/} seguendo l'ordine alfabetico dei nomi dei file. Per questo motivo, i file sono stati nominati con un prefisso numerico:

\begin{enumerate}
    \item \textbf{\texttt{01\_init.sql}}: Contiene i comandi \texttt{DROP DATABASE IF EXISTS} e \texttt{CREATE DATABASE} per garantire un ambiente pulito, seguiti dalla creazione delle tabelle con i relativi vincoli di integrità.
    \item \textbf{\texttt{02\_dump.sql}}: Contiene gli \texttt{INSERT} iniziali per popolare la tabella \texttt{corsi}, gli argomenti di test e gli utenti predefiniti.
\end{enumerate}

\section{Script di Automazione}
Per semplificare la gestione dell'ambiente, sono stati predisposti tre script \texttt{.sh} nella root del progetto.

\subsection{Avvio del Progetto (\texttt{run\_project.sh})}
Utilizzare questo script per avviare l'intero stack. Lo script richiede all'utente se desidera effettuare un reset dei volumi per ricaricare il database da zero.
\begin{verbatim}
./run_project.sh
\end{verbatim}

\subsection{Arresto del Progetto (\texttt{stop\_project.sh})}
Ferma i container in esecuzione senza eliminare i dati salvati nel database.
\begin{verbatim}
./stop_project.sh
\end{verbatim}

\subsection{Reset Totale (\texttt{reset\_all.sh})}
Elimina container, volumi e immagini, ricostruendo l'ambiente da uno stato pulito. Utile in caso di modifiche strutturali al file \texttt{init.sql}.
\begin{verbatim}
./reset_all.sh
\end{verbatim}

\section{Procedura di Installazione Standard}
\begin{enumerate}
    \item Clonare il repository o estrarre i file del progetto.
    \item Aprire il terminale nella cartella principale.
    \item Fornire i permessi di esecuzione agli script:
    \begin{verbatim}
chmod +x *.sh
    \end{verbatim}
    \item Eseguire lo script di avvio:
    \begin{verbatim}
./run_project.sh
    \end{verbatim}
    \item Almeno al primo avvio alla domanda su reset rispondere S
    \begin{verbatim}
Vuoi resettare il database e lo storage? (Perderai i dati correnti) [S/N]
    \end{verbatim}
    \item Attendere circa 10-15 secondi per consentire a MySQL di completare l'inizializzazione.
    \item Accedere all'app tramite browser all'indirizzo:
        \begin{verbatim}
http://localhost:8080
        \end{verbatim}
\end{enumerate}

\section{Credenziali di Test}
Per verificare le funzionalità del sistema (Login, consultazione e permessi), sono stati predisposti i seguenti account nel file di dump:


\begin{table}[h!]
\centering
\caption{Credenziali di test pre-caricate nel sistema.}
\label{tab:credenziali}
\resizebox{\textwidth}{!}{% <-- Inizio ridimensionamento
\begin{tabular}{|l|l|l|l|}
\hline
\textbf{Tipo Utente} & \textbf{Email} & \textbf{Password} & \textbf{Ruolo} \\ \hline
Amministratore & \texttt{admin@unipr.it} & \texttt{Vector!} & \texttt{amministratore} \\ \hline
Studente & \texttt{studente\_test1@studenti.unipr.it} & \texttt{Prova1} & \texttt{studente} \\ \hline
Studente & \texttt{studente\_test2@studenti.unipr.it} & \texttt{Prova2} & \texttt{studente} \\ \hline
Studente & \texttt{studente\_test3@studenti.unipr.it} & \texttt{Prova3} & \texttt{studente} \\ \hline
\end{tabular}% <-- Fine ridimensionamento
}
\end{table}


\section{Risoluzione dei Problemi}
Se le tabelle non dovessero risultare presenti (\textit{Table not found}), eseguire il comando \texttt{./reset\_all.sh}. Questo forzerà Docker a cancellare il volume dati esistente e a rieseguire gli script di inizializzazione \texttt{01\_init.sql} e \texttt{02\_dump.sql}.

